\documentclass{article}
\usepackage{amsmath,amssymb,hyperref,graphicx}
\title{Discovery Engine: running draft}
\date{\today}
\begin{document}
\maketitle


\section{Abstract}

We ask whether small transformers trained to predict simulated orbits ($G = M = 1$) encode the physical law, and when such structure forms. Linear probes target force $m_2/r^2$, per-snapshot acceleration \texttt{accel\_step} $(F/m_2)\,dt^2$, distance and energy. A result counts only if it beats several random-model controls. The lab has 45 replicated findings (32 supported, 13 refuted) from 107 evidence runs, 4 of them cut short by the budget stop. \texttt{accel\_step} is decodable in every setting. At $N = 64$, $e \le 0.5$, it forms at steps 85–140, which is 40–90 steps before the model beats constant-velocity extrapolation [\texttt{f-20261012-002713-8d93}, \texttt{f-20261012-005950-42a7}]. Its headroom-normalised margin stays at about 0.63–0.79 from $N = 16$ to $N = 256$ [\texttt{f-20261013-000329-a86b}, \texttt{f-20261013-075824-35f0}, \texttt{f-20261013-063603-95b5}]. Force is weaker. Its margin falls from 0.216–0.255 at $N = 64$ to 0.079–0.143 at $N = 128$, and the loss is in trained $R^2$, not in the controls. At $N = 256$ it falls to 0.026–0.045 [\texttt{f-20261013-075855-f75f}, \texttt{f-20261013-075824-35f0}, \texttt{f-20261013-075749-3d2c}]. The main negative result is that symbolic regression on probe directions returns $r^{-1}$ forms such as $a^2/r - 0.55a$, never the true $a^3/r^2$ [\texttt{f-20261011-231118-7a37}, \texttt{f-20261011-223027-8fb3}]. The fitter has not yet been tested on the true target.

\section{Methods}

\subsection{Data generation}

All orbits are simulated with REBOUND in units where $G = 1$ and the primary mass is $M = 1$. The gravitational force on body 2 is therefore $|F| = m_2/r^2$, where $m_2$ is the mass ratio. Varying $M$ would add nothing. Positions are sampled at fixed phase intervals, so absolute mass and time scale cannot be identified from them.

Each orbit is sampled at a fixed number of snapshots per orbit, $N$ (\texttt{steps\_per\_orbit}). The time between snapshots is therefore $\Delta t = P/N$, which grows with orbit size ($\Delta t^2 \propto a^3$). The model never sees physical time.

Three-body systems are hierarchical: the third body sits at 3.5–6× the inner semi-major axis. They are used only as an evaluation set for models trained on two-body systems.

The configuration schema adds these fields to the specification:

\begin{itemize}
\item \texttt{data.eccentricity\_range}
\item \texttt{data.data\_seed}
\item \texttt{probes.control\_margin}
\item \texttt{probes.patching}
\item \texttt{probes.symbolic}
\item \texttt{eval\_sets}, for out-of-distribution evaluation
\end{itemize}

Unknown fields are rejected. The training ranges differ between the cited experiments:

\begin{itemize}
\item Some runs train on the default range, which reaches $e \leq 0.5$ [\texttt{f-20261011-223027-8fb3}].
\item Some runs widen the mass ratio to $m_2 \in [0.01, 1]$ [\texttt{f-20261011-203628-5cd9}].
\item Some runs widen eccentricity to $e \in [0, 0.8]$ [\texttt{f-20261011-211551-d718}].
\item Some runs widen both, to $e \in [0, 0.9]$ and $m_2 \in [0.01, 1]$ [\texttt{f-20261011-223027-8fb3}]. Both of these runs share data seed 0.
\end{itemize}

The held-out evaluation sets are:

\begin{itemize}
\item low mass ratio, $m_2 \in [0.01, 0.1]$
\item high eccentricity, $e \in [0.5, 0.8]$
\item three-body systems [\texttt{f-20261011-203723-b7de}]
\item held-out low-$e$ and high-$e$ splits, used for the wide-eccentricity runs [\texttt{f-20261011-211551-d718}]
\end{itemize}

\subsection{Tokens, model and prediction target}

Each token is one time step. It holds the $xy$ positions of up to three bodies (6 numbers) in the centre-of-mass frame, with absent bodies set to zero. Because every system uses the same token format, a model trained on two bodies can be evaluated on three-body sets.

The model predicts the normalised displacement to the next step. All evidence runs use one transformer architecture: $L = 4$ layers, width $d = 64$, and $N = 64$ snapshots per orbit [\texttt{f-20261011-211551-d718}].

Prediction quality is reported as \texttt{val\_mse\_ratio}. This is the model's MSE divided by the MSE of constant-velocity extrapolation, $\hat{x}_{t+1} = 2x_t - x_{t-1}$. Both are computed in physical position units over steps $t \geq 1$ on the same held-out systems. A value below 1 means the model beats the baseline.

\subsection{Training}

The lab's default first-experiment (E1) configuration trains for 5k steps at batch size 128 and learning rate $10^{-3}$. Training runs on a laptop GPU (MPS) and is kept CPU-friendly.

Checkpoint spacing varies between experiments. One formation experiment probes every 250 steps [\texttt{f-20261011-201533-0653}]. A dense-checkpoint experiment probes every 25 steps early in training [\texttt{f-20261011-203659-4e2f}].

Every finding is replicated on two training seeds.

\textbf{Budget hard stop.} A runaway guard stops a run once it reaches 2× its seed's share of \texttt{budget.max\_compute\_hours}. The clock starts at the first training step, so dataset generation does not count. A stopped run is checkpointed and probed at its stop step. It records \texttt{steps\_completed} and \texttt{steps\_planned}, and its \texttt{completion} is the fraction of planned steps it trained. Runs below 80\% completion do not count toward replication.

Two of the 12 evidence runs were cut short by an earlier version of this guard. That version stopped at 1.2× the budget share and counted dataset generation against it. Both runs are in [\texttt{f-20261011-211551-d718}]: seed 9 stopped at completion 0.867 (step 2640 of 3000) and seed 8 at 0.9 (step 2714). Under that version the probes were taken at the last regular checkpoint, step 2500, not at the stop step. Their \texttt{val\_mse\_ratio} was still falling when training stopped.

\subsection{Probe targets}

Linear probes read five quantities from the residual stream at every layer:

\begin{itemize}
\item \texttt{distance}
\item relative \texttt{velocity}
\item \texttt{energy}, which is constant per system
\item physical \texttt{force}, $|F| = m_2/r^2$
\item the per-snapshot acceleration \texttt{accel\_step}:
\end{itemize}

\[|a_{\text{step}}| = \frac{F}{m_2}\,\Delta t^2 = \left(\frac{2\pi}{N}\right)^2 \frac{a^3}{(1+m_2)\,r^2}.\]

Next-step prediction needs \texttt{accel\_step}, not the physical force. A weak \texttt{force} probe on its own could therefore mean the model learned a per-snapshot version of gravity, not that it learned no gravity. For this reason \texttt{force} and \texttt{accel\_step} are always compared.

A unit test checks two properties of \texttt{accel\_step}:

\begin{itemize}
\item It equals the second difference of positions, to within 1\%.
\item It is not proportional to the physical force across orbit sizes.
\end{itemize}

\texttt{velocity} probe results do not appear in any cited finding.

\subsection{Probe protocol and random-model controls}

Probes run on 300 held-out systems, at every other time step from $t = 8$. Systems are split 70/30 into probe-train and probe-test sets, so rows from the same orbit never appear on both sides.

Each probe is a ridge regression (RidgeCV over $\alpha \in [10^{-3}, 10^{4}]$) on standardised inputs and outputs. $R^2$ is averaged across output dimensions.

Every probe has a control: the same probe fitted on a randomly initialised model of the same architecture. There is one control per layer, because a random model does not change across checkpoints. A probe result counts only if

\[R^2 - R^2_{\text{control}} \geq 0.1 \quad (\texttt{probes.control\_margin}).\]

For each target, the lab reports:

\begin{itemize}
\item the best layer
\item its $R^2$ and control $R^2$
\item \texttt{formed\_at\_step}, the first checkpoint at which the margin reaches the threshold (null if it never does)
\end{itemize}

The control $R^2$ for \texttt{accel\_step} is high under the default ranges (0.71–0.75). This caps the margin \texttt{accel\_step} can reach. One finding therefore also reports a post hoc, headroom-normalised score, $(R^2 - R^2_{\text{control}})/(1 - R^2_{\text{control}})$, but its pre-registered criterion uses the raw margin [\texttt{f-20261011-203628-5cd9}].

Probes on the out-of-distribution evaluation sets have no random controls. Their $R^2$ values therefore cannot be read as margins [\texttt{f-20261011-203723-b7de}, \texttt{f-20261011-211551-d718}].

Only the final checkpoint's activations are saved to disk, as float16. Activations for earlier checkpoints are recomputed from the saved weights when needed.

\subsection{Symbolic regression}

When \texttt{probes.symbolic} is enabled, PySR fits a closed-form expression to the best-layer probe direction for \texttt{force} and for \texttt{accel\_step}. The \texttt{accel\_step} fit uses the inputs $r$, $m_2$ and $a$. Each fit reports:

\begin{itemize}
\item the equation
\item whether the true law was recovered (\texttt{law\_recovered})
\item \texttt{fit\_error\_vs\_true\_law}
\end{itemize}

Symbolic fits were run at the final checkpoint only, and only in the runs behind [\texttt{f-20261011-223027-8fb3}]. These are the only evidence runs with fitted equations.

\subsection{Activation patching}

The pipeline also supports activation patching. Heads are ranked by a force probe on each head's contribution to the residual stream. The top three heads, plus the lowest-ranked head as a control, are each patched with that head's output on a different system, and the relative increase in next-step MSE is recorded. None of the cited findings report patching results.

\section{Results}

The lab has 45 citable findings, and all 45 are replicated. Of these, 32 support their hypothesis and 13 refute it. Together they rest on 107 evidence runs. The budget hard stop cut 4 of those runs short, and we say so wherever their numbers appear. Unless stated otherwise, the model has 4 layers and $d_\text{model} = 64$. "Margin" means probe $R^2$ minus the best random-control $R^2$. A probe beats its control when the margin is at least the required margin, $\max(0.1, \text{control spread})$. Where controls differ a lot between targets, we also report the headroom-normalised margin, $(R^2 - \text{control})/(1 - \text{control})$.

\subsection{Prediction quality}

Across the configurations tested, the models beat constant-velocity extrapolation on held-out systems. Completed $N = 64$ runs at $e \le 0.5$ reach \texttt{val\_mse\_ratio} of about 0.011–0.016 after 3000–4000 steps [\texttt{f-20261011-201533-0653}, \texttt{f-20261013-075855-f75f}]. With 12000 steps this falls to about 0.003 [\texttt{f-20261013-051738-2a37}]. A smaller model (2 layers, $d_\text{model} = 32$, 1000 systems) reaches 0.064 and 0.079 [\texttt{f-20261013-000220-ab2d}]. The exception is $N = 256$ at $e \le 0.85$ with only 600 steps, which never beats the baseline: its final ratio is 2.03–2.37 [\texttt{f-20261012-034134-5e82}].

\subsection{What forms first, and when}

\textbf{Geometry forms first, then \texttt{accel\_step}, and force forms late or not at all.} This ordering holds in every timing experiment.

\begin{itemize}
\item \textbf{Distance and energy} cross their controls within steps 5–50 [\texttt{f-20261011-203659-4e2f}, \texttt{f-20261012-002713-8d93}, \texttt{f-20261012-005950-42a7}, \texttt{f-20261012-044142-d490}].
\item \textbf{\texttt{accel\_step}} forms around steps 60–140. With dense checkpoints at $N = 64$ and $e \le 0.5$, it first holds the 0.1 margin at:
\item steps 100 in two seeds, with 25-step checkpoints [\texttt{f-20261011-203659-4e2f}];
\item steps 140, 140 and 100 [\texttt{f-20261012-002713-8d93}];
\item steps 85–90 [\texttt{f-20261012-005950-42a7}];
\item steps 90, 100 and 100 [\texttt{f-20261012-044142-d490}];
\item steps 65–90 with 5-step checkpoints [\texttt{f-20261012-054929-0067}].
\item \textbf{Before or after beating the baseline.} In these runs, \texttt{accel\_step} forms 40–90 steps before \texttt{val\_mse\_ratio} first falls below 1. At that point the ratio is still about 1.7–2.1 [\texttt{f-20261012-002713-8d93}, \texttt{f-20261012-044142-d490}].
\item \textbf{Force} fails its control at the end of these 300-step runs, with final margins of 0.000–0.080 [\texttt{f-20261012-002713-8d93}, \texttt{f-20261012-044142-d490}].
\end{itemize}

!\href{figure:formation/run-20261012-003058-0ea4}{Probe formation, N = 64, seed 47}

\textbf{Several pre-registered refinements failed.}

\begin{itemize}
\item \textbf{A stricter bar.} With best-of-5 controls, the \texttt{accel\_step} margin levelled off at 0.134–0.135 in 2 of 3 seeds, below a 0.15 bar [\texttt{f-20261012-005950-42a7}].
\item \textbf{The early force spike is real.} A force margin of 0.166–0.189 appears at step 10, then decays to about 0.00–0.02 by steps 40–50. It survives best-of-5 controls (refuted) [\texttt{f-20261012-005950-42a7}].
\item \textbf{A narrow formation window.} Only 2 of 5 seeds landed in the pre-registered $N = 64$ window (refuted) [\texttt{f-20261012-054929-0067}]. In the same experiment, force barely cleared its control in 2 of 5 seeds (margins 0.103 and 0.105).
\item \textbf{Learning rate.} Halving the learning rate delayed formation by only 10–30 steps, not the roughly 2× that proportional scaling predicts (refuted) [\texttt{f-20261012-044258-ab88}].
\end{itemize}

\textbf{"Beating the baseline" is not a stable reference point across $N$.} How strong the constant-velocity baseline is depends on $N$, so the ordering of the two events changes while the formation step does not.

\begin{itemize}
\item \textbf{At $N = 16$}, \texttt{val\_mse\_ratio} crosses 1 at step 20. That is well before \texttt{accel\_step} forms at steps 75–90, and those formation steps are within 15 steps of the seed-matched $N = 64$ runs [\texttt{f-20261012-015955-fc2b}].
\item \textbf{At $e \le 0.85$}, \texttt{accel\_step} forms at steps 60–70 at $N = 64$ [\texttt{f-20261012-034018-5526}]. It forms at steps 60–70 at $N = 256$ too, even though those runs never beat the baseline [\texttt{f-20261012-034134-5e82}].
\end{itemize}

So formation follows a training-step window, not relative loss reduction [\texttt{f-20261012-034134-5e82}].

\textbf{When training runs longer, force forms persistently, but late.}

\begin{itemize}
\item \textbf{Coarse checkpoints.} Even with 250-step checkpoints and a single control, gravity-related probes already beat their controls at step 250 [\texttt{f-20261011-201533-0653}].
\item \textbf{Larger model.} With 4 layers and $d_\text{model} = 128$ on data seed 1, force forms at steps 400 and 350, against step 100 for \texttt{accel\_step}. It keeps strengthening up to 5000 steps, reaching margins of 0.187 and 0.208 [\texttt{f-20261013-020028-8792}].
\item \textbf{Longer default runs.} In 1500-step runs on data seed 1, force first holds the bar at steps 400–600. This happens only after its best layer moves to layer 1, the layer where distance is best decoded [\texttt{f-20261013-025812-b9ee}].
\end{itemize}

\subsection{\texttt{accel\_step} versus force}

\textbf{Which probe "wins" depends on the metric.} In complete $N = 64$, $e \le 0.5$ runs:

\begin{itemize}
\item \textbf{Raw $R^2$:} \texttt{accel\_step} is higher (about 0.90–0.94, against 0.68–0.74 for force) [\texttt{f-20261013-075855-f75f}, \texttt{f-20261013-054512-1b08}].
\item \textbf{Raw margin:} force often leads, because its control is much lower (about 0.45–0.53, against 0.69–0.79 for \texttt{accel\_step}) [\texttt{f-20261013-075855-f75f}, \texttt{f-20261013-054512-1b08}].
\item \textbf{Headroom-normalised margin:} \texttt{accel\_step} leads in every seed. It is 0.68–0.74, against 0.43–0.48 for force [\texttt{f-20261013-075855-f75f}, \texttt{f-20261013-054512-1b08}].
\end{itemize}

\textbf{Two related checks:}

\begin{itemize}
\item \textbf{Wider $m_2$ range.} Widening $m_2$ to $[0.01, 1]$ did not put force behind \texttt{accel\_step} on raw margin (refuted) [\texttt{f-20261011-203628-5cd9}].
\item \textbf{Smaller model} (2 layers, $d_\text{model} = 32$). Force has the larger raw margin here too (0.221 and 0.255, against 0.146 and 0.128). On normalised margin the ordering splits between the two seeds [\texttt{f-20261013-000220-ab2d}].
\end{itemize}

\textbf{Where the force signal lives.} Force usually peaks at layer 1, the layer where distance is best decoded. No force probe with distance regressed out has been run. The force signal may therefore partly be a nonlinear reading of distance [\texttt{f-20261013-015914-3f49}, \texttt{f-20261013-025812-b9ee}, \texttt{f-20261013-051738-2a37}].

!\href{figure:heatmap/run-20261013-074422-d66a/accel_step}{accel\_step probe heatmap, N = 64, seed 24}

\subsection{Sampling density: force decodability falls with $N$, \texttt{accel\_step} does not}

Final force margins, after 3000 steps unless noted:

| | $N = 16$ | $N = 64$ | $N = 128$ | $N = 256$ |
|---|---|---|---|---|
| Data seed 0, seeds 20/21 | 0.236 / 0.275 [\texttt{f-20261013-000329-a86b}] | 0.215 / 0.261 [\texttt{f-20261013-015914-3f49}] | 0.150 / 0.131 [\texttt{f-20261013-051807-f3f1}] | 0.113 / 0.003 (rerun) [\texttt{f-20261013-050513-2c7a}] |
| Data seed 0, seeds 22/23 | — | 0.247 / 0.234 [\texttt{f-20261013-015914-3f49}] | 0.108 / 0.094 [\texttt{f-20261013-051807-f3f1}] | 0.040 / 0.033 [\texttt{f-20261013-063603-95b5}] |
| Data seed 0, seeds 24/25/26 | — | 0.228 / 0.255 / 0.216 [\texttt{f-20261013-075855-f75f}] | 0.128 / 0.143 / 0.079 [\texttt{f-20261013-075824-35f0}] | 0.026 / 0.045 / 0.029 [\texttt{f-20261013-075749-3d2c}] |
| Data seed 1, seeds 20/21 | — | 0.209 / 0.209 [\texttt{f-20261013-054512-1b08}] | 0.137 / 0.145 [\texttt{f-20261013-054542-23e1}] | 0.072 / 0.039 [\texttt{f-20261013-063631-c4ed}] |
| Data seed 1, seeds 22/23 | — | 0.268 / 0.275 (trajectory) [\texttt{f-20261013-224908-8b46}] | 0.173 / 0.163 [\texttt{f-20261013-224908-8b46}] | — |
| $m_2 = 0.5$, data seed 0, seeds 20/21 | — | 0.133 / 0.195 [\texttt{f-20261013-234439-d851}] | 0.063 / 0.101 [\texttt{f-20261013-234552-f8e8}] | — |
| $m_2 = 0.5$, data seed 0, seeds 22/23 | — | 0.193 / 0.228 (paired arm) [\texttt{f-20261014-004426-68d6}] | 0.062 / 0.080 [\texttt{f-20261014-004426-68d6}]; 6000 steps: 0.088 / 0.145 [\texttt{f-20261014-025914-c84b}] | — |
| $m_2 = 0.5$, data seed 1, seeds 20/21 | — | 0.150 / 0.175 [\texttt{f-20261014-015644-a80b}] | 0.092 / 0.0999 [\texttt{f-20261014-015607-930a}] | — |

\textbf{$N = 16$ and $N = 64$ give the same force level.} The hypothesis that force declines gradually from $N = 16$ was refuted [\texttt{f-20261013-015914-3f49}].

\textbf{From $N = 64$ to $N = 128$, force falls in every paired seed. The loss is in trained $R^2$, not in the controls.}

\begin{itemize}
\item \textbf{Data seed 0.} For seeds 24–26, trained force $R^2$ fell by 0.098–0.128, while controls moved by only +0.002 to +0.009 [\texttt{f-20261013-075824-35f0}].
\item \textbf{Data seed 1.} For seeds 22/23, trained $R^2$ fell by 0.095–0.097, while controls moved by −0.002 to +0.017 [\texttt{f-20261013-224908-8b46}]. For seeds 20/21 the drop is smaller, 0.064–0.072, with flat controls [\texttt{f-20261013-054542-23e1}].
\item \textbf{Fixed $m_2 = 0.5$.} The drop persists, so the hypothesis that fixing $m_2$ removes it was refuted. Trained $R^2$ fell by:
\item 0.070 / 0.105 on data seed 0 [\texttt{f-20261013-234552-f8e8}];
\item about 0.11 / 0.14 on data seed 0, seeds 22/23 [\texttt{f-20261014-004426-68d6}];
\item about 0.065 / 0.088 on data seed 1, with same-layer controls moving by only −0.007 / −0.013 [\texttt{f-20261014-015607-930a}].
\item \textbf{Force fails its control at $N = 128$ with fixed $m_2$} in 5 of 6 seeds after 3000 steps [\texttt{f-20261013-234552-f8e8}, \texttt{f-20261014-004426-68d6}, \texttt{f-20261014-015607-930a}].
\item \textbf{Interpretation.} What is lost at $N = 128$ is the readout of $1/r^2$, not the $m_2$ factor [\texttt{f-20261013-234552-f8e8}].
\end{itemize}

\textbf{Longer training narrows the $N = 64$ to $N = 128$ gap but has not closed it.}

\begin{itemize}
\item \textbf{12000 steps, wide $m_2$.} $N = 64$ stays 0.062 / 0.069 above $N = 128$ [\texttt{f-20261013-051738-2a37}].
\item \textbf{6000 steps, fixed $m_2$.} We trained $N = 128$ to 6000 steps for seeds 22/23. Trained force $R^2$ reached 0.681 / 0.696, still 0.086 / 0.074 below the 3000-step $N = 64$ values [\texttt{f-20261014-025914-c84b}]. This met the pre-registered ceiling criterion. But between steps 3000 and 6000 the gap shrank from 0.098 to 0.086 (seed 22) and from 0.123 to 0.074 (seed 23). Seed 23 is still gaining about 0.005–0.006 per 500 steps. This comparison also sets 6000 $N = 128$ steps against 3000 $N = 64$ steps, so a fixed ceiling has not been shown.
\end{itemize}

\textbf{From $N = 128$ to $N = 256$, what drives the drop depends on the dataset.}

\begin{itemize}
\item \textbf{Data seed 1.} The drop is entirely a control effect. Trained $R^2$ stays at about 0.66–0.67 while force controls rise to 0.60–0.63 [\texttt{f-20261013-063631-c4ed}].
\item \textbf{Data seed 0.} Trained $R^2$ falls as well. Controls explain only about 35–50\% of the margin loss, so the controls-only hypothesis was refuted [\texttt{f-20261013-075749-3d2c}].
\end{itemize}

\textbf{On data seed 0, the $N = 256$ shortfall is not just slower convergence.}

\begin{itemize}
\item \textbf{1000 extra steps} added only 0.011–0.014 of trained force $R^2$ (refuted) [\texttt{f-20261013-083108-2efd}].
\item \textbf{Matched on \texttt{val\_mse\_ratio}}, $N = 128$ still decodes force better: 0.641 vs 0.557, and 0.630 vs 0.569 [\texttt{f-20261013-093210-2e03}].
\end{itemize}

!\href{figure:heatmap/run-20261013-071946-8084/force}{Force probe heatmap, N = 256, seed 24}

\textbf{The normalised \texttt{accel\_step} margin stays roughly flat across $N$:}

| $N$ | Normalised \texttt{accel\_step} margin | Sources |
|---|---|---|
| 16 | 0.68–0.69 | [\texttt{f-20261013-000329-a86b}] |
| 64 | 0.68–0.79 | [\texttt{f-20261013-015914-3f49}, \texttt{f-20261013-075855-f75f}, \texttt{f-20261013-054512-1b08}, \texttt{f-20261013-234439-d851}, \texttt{f-20261014-015644-a80b}] |
| 64, 12000 steps | 0.76 / 0.65 | [\texttt{f-20261013-051738-2a37}] |
| 128 | 0.63–0.78 | [\texttt{f-20261013-051807-f3f1}, \texttt{f-20261013-075824-35f0}, \texttt{f-20261013-054542-23e1}, \texttt{f-20261013-224908-8b46}, \texttt{f-20261013-234552-f8e8}, \texttt{f-20261014-004426-68d6}, \texttt{f-20261014-015607-930a}, \texttt{f-20261014-025914-c84b}] |
| 256 | 0.66–0.75 | [\texttt{f-20261013-063603-95b5}, \texttt{f-20261013-063631-c4ed}, \texttt{f-20261013-075749-3d2c}, \texttt{f-20261013-083108-2efd}] |

At $N = 16$, the hypothesis that coarser sampling would raise this margin was refuted [\texttt{f-20261013-000329-a86b}].

\subsection{Fixing the mass ratio}

\textbf{Fixing $m_2 = 0.5$ does not lift force to the hypothesised $R^2 \ge 0.85$ (refuted).} In the first test, force reached 0.750 / 0.732 while \texttt{accel\_step} reached 0.924 / 0.946 [\texttt{f-20261011-223000-b5f9}]. The budget hard stop cut both of those runs short:

\begin{itemize}
\item seed 14 stopped at step 1500 (completion 0.55), below the 80\% needed to count toward replication;
\item seed 15 stopped at step 2500 (completion 0.917).
\end{itemize}

\textbf{In complete 3000-step $N = 64$ runs, trained force $R^2$ is:}

\begin{itemize}
\item 0.763 / 0.789 on data seed 0 [\texttt{f-20261013-234439-d851}];
\item 0.795 / 0.790 on data seed 1 [\texttt{f-20261014-015644-a80b}].
\end{itemize}

These are higher than in wide-$m_2$ runs, but the margins are smaller. Fixing $m_2$ raises the force control to about 0.59–0.64 [\texttt{f-20261013-234439-d851}, \texttt{f-20261014-015644-a80b}].

\subsection{Eccentricity sets the control level}

Widening the training range of $e$ lowers both the \texttt{accel\_step} control and its trained $R^2$. All rows are $N = 64$:

| $e$ range | \texttt{accel\_step} control $R^2$ | Trained $R^2$ | Normalised margin |
|---|---|---|---|
| $\le 0.5$ [\texttt{f-20261012-044142-d490}] | 0.724–0.746 | 0.890–0.902 | — |
| $\le 0.7$ [\texttt{f-20261012-055051-b178}] | 0.543–0.559 | 0.80–0.81 | ≈0.55 |
| $\le 0.75$ [\texttt{f-20261013-005923-526c}] | 0.463–0.487 | 0.746–0.753 | 0.51–0.54 |
| $\le 0.8$ [\texttt{f-20261013-000359-37e2}] | 0.372–0.396 | 0.670–0.672 | 0.455–0.475 |
| $\le 0.85$ [\texttt{f-20261012-034018-5526}] | 0.248–0.333 | 0.507–0.594 | 0.34–0.40 |

\begin{itemize}
\item \textbf{The decline steepens above $e \le 0.75$} [\texttt{f-20261013-005923-526c}].
\item \textbf{Force is weaker at every ceiling.} Its normalised margins are 0.11–0.22, and it forms late, at steps 250–590 [\texttt{f-20261012-034018-5526}, \texttt{f-20261012-055051-b178}, \texttt{f-20261013-000359-37e2}, \texttt{f-20261013-005923-526c}].
\item \textbf{Larger model at $e \le 0.85$} (6 layers, $d_\text{model} = 128$). Linear \texttt{accel\_step} $R^2$ levels off at about 0.70–0.74, gaining only 0.022–0.025 after step 2000. Over the same steps, \texttt{val\_mse\_ratio} keeps falling 5.5–6.8× [\texttt{f-20261012-030233-ab28}]. So linear decodability stops tracking prediction quality.
\end{itemize}

\subsection{Out-of-distribution transfer}

\textbf{Trained on $e \le 0.5$} [\texttt{f-20261011-203723-b7de}]:

\begin{itemize}
\item \textbf{Low mass ratios} ($m_2 \in [0.01, 0.1]$). Error is about 3.6× the in-distribution error, and force $R^2$ holds at 0.75–0.77.
\item \textbf{High eccentricity} ($e \in [0.5, 0.8]$). Error is 16–18× the in-distribution error, which fails the pre-registered 10× criterion (refuted). The models still beat the baseline. Force $R^2$ falls to 0.40 / 0.34.
\item \textbf{Three-body sets.} Ratios of 72 and 70 are exploratory only.
\end{itemize}

\textbf{Trained on $e \in [0, 0.8]$} [\texttt{f-20261011-211551-d718}]:

\begin{itemize}
\item \textbf{The high-$e$ gap closes.} High-$e$ error falls to 1.01–1.05× the in-distribution error.
\item \textbf{These runs were cut short.} The budget stop ended both at step 2500 (completion 0.867 and 0.9), so their probes describe the step-2500 weights.
\item \textbf{The probes still differ between regimes.} With prediction error matched, force $R^2$ on held-out high-$e$ systems is 0.366–0.377, against 0.729–0.802 on low-$e$ systems. These out-of-distribution probes have no random controls.
\end{itemize}

\subsection{Symbolic regression has not recovered $r^{-2}$}

\begin{itemize}
\item \textbf{Strong probes} (\texttt{accel\_step} $R^2$ 0.831 / 0.881). The fitted direction has the form $a^2/r - c\,a$ with $c \approx 0.55$, not the true law $a^3/r^2$. Fit errors are 0.245 / 0.201. Force fits as about $0.8/r$ [\texttt{f-20261011-231118-7a37}].
\item \textbf{Weaker probes, training at $e \le 0.9$} (\texttt{accel\_step} $R^2$ 0.40–0.50). These also fit as $r^{-1}$; the expected $r^{-2}$ did not appear (refuted) [\texttt{f-20261011-223027-8fb3}].
\item \textbf{Possible fitter artefact.} The fitter has not been run on the true \texttt{accel\_step} target. We cannot yet rule out that the fitter itself produces the $r^{-1}$ form.
\end{itemize}

\subsection{Reproducibility}

\textbf{The original $N = 256$ reference runs did not reproduce.} An exact rerun gave force margins of 0.113 / 0.003. At step 3000 it agreed with 6000- and 12000-step runs to within about 0.025. That rules out the learning-rate schedule horizon as the cause (refuted) [\texttt{f-20261013-050513-2c7a}].

\textbf{The current pipeline is deterministic for a fixed seed.} A second identical rerun matched the first to 3 decimals at every checkpoint [\texttt{f-20261013-054603-d6cc}]. Two things follow:

\begin{itemize}
\item the 0.11 force gap between seeds 20 and 21 at $N = 256$ is genuine seed variance;
\item the cause of the original discrepancy is still unexplained.
\end{itemize}

\subsection{Null and negative results}

The 13 refuted hypotheses:

\begin{enumerate}
\item \textbf{A wider $m_2$ range would put force behind \texttt{accel\_step}.} It did not [\texttt{f-20261011-203628-5cd9}].
\item \textbf{A model trained on $e \le 0.5$ would transfer to high $e$ within 10×.} The error was 16–18× [\texttt{f-20261011-203723-b7de}].
\item \textbf{Fixing $m_2$ would lift force $R^2$ to 0.85.} It did not [\texttt{f-20261011-223000-b5f9}].
\item \textbf{Wide-$e$ training would yield an $r^{-2}$ fit.} The fits were $r^{-1}$ [\texttt{f-20261011-223027-8fb3}].
\item \textbf{\texttt{accel\_step} would clear a 0.15 bar, and the early force spike would vanish under best-of-5 controls.} Neither happened [\texttt{f-20261012-005950-42a7}].
\item \textbf{Formation would scale in proportion to the learning rate.} The shift was sub-proportional [\texttt{f-20261012-044258-ab88}].
\item \textbf{$N = 64$ formation would fall in a narrow window.} Only 2 of 5 seeds did [\texttt{f-20261012-054929-0067}].
\item \textbf{$N = 16$ would raise the normalised \texttt{accel\_step} margin.} It was level with $N = 64$ [\texttt{f-20261013-000329-a86b}].
\item \textbf{Force would decline gradually from $N = 16$.} $N = 64$ matched $N = 16$ [\texttt{f-20261013-015914-3f49}].
\item \textbf{An exact rerun would reproduce the original $N = 256$ margins.} It did not [\texttt{f-20261013-050513-2c7a}].
\item \textbf{The $N = 128 \to 256$ force drop is all a control effect.} On data seed 0, trained $R^2$ also fell [\texttt{f-20261013-075749-3d2c}].
\item \textbf{1000 extra steps would close the $N = 256$ gap.} They did not [\texttt{f-20261013-083108-2efd}].
\item \textbf{Fixing $m_2$ would remove the $N = 64 \to 128$ drop.} The drop persisted [\texttt{f-20261013-234552-f8e8}].
\end{enumerate}

Several supported findings are also nulls for force, in that force fails its control:

\begin{itemize}
\item short $N = 64$ runs [\texttt{f-20261012-002713-8d93}, \texttt{f-20261012-044142-d490}];
\item $N = 256$ in most seeds [\texttt{f-20261012-034134-5e82}, \texttt{f-20261013-063603-95b5}, \texttt{f-20261013-063631-c4ed}];
\item $m_2$-fixed runs at $N = 128$ [\texttt{f-20261014-004426-68d6}, \texttt{f-20261014-015607-930a}], and seed 22 even after 6000 steps [\texttt{f-20261014-025914-c84b}].
\end{itemize}

\subsection{Not yet tested, and limits}

The lab has not yet reported:

\begin{itemize}
\item \texttt{velocity} probe results;
\item activation patching, or any other causal test of whether the decoded codes are used;
\item a force probe with \texttt{distance} and \texttt{accel\_step} regressed out;
\item a random-control pool shared across $N$;
\item nonlinear or log-target probes;
\item symbolic regression on the true \texttt{accel\_step} target;
\item matched-loss and matched-horizon comparisons for data seed 1 and for the $m_2$-fixed arms.
\end{itemize}

Limits that apply throughout:

\begin{itemize}
\item \textbf{Probes.} All probes are linear.
\item \textbf{Single controls.} The earliest findings used a single random control [\texttt{f-20261011-201533-0653}, \texttt{f-20261011-203628-5cd9}, \texttt{f-20261011-203659-4e2f}].
\item \textbf{Formation steps.} \texttt{formed\_at\_step} sometimes lags the trajectory by one checkpoint [\texttt{f-20261012-054929-0067}, \texttt{f-20261013-000329-a86b}, \texttt{f-20261013-054512-1b08}].
\item \textbf{Rounded values.} Several paired comparisons rely on 3-decimal trajectory values whose $R^2$ field may be a maximum over layers [\texttt{f-20261013-075749-3d2c}, \texttt{f-20261014-004426-68d6}, \texttt{f-20261014-015607-930a}].
\item \textbf{Training length.} Force was still rising at the end of most runs.
\item \textbf{Datasets.} Only two data seeds have been tested.
\item \textbf{Architecture.} Most results come from one architecture (4 layers, $d_\text{model} = 64$).
\end{itemize}

\end{document}
